\section{Impacto Académico}
\LaTeX procesa alrededor de 85\% de papers en arXiv.org y es requerido por journals como IEEE Transactions, Nature y Springer \cite{ieee_requirements,arxiv_stats}. Su importancia radica en la capacidad única para manejar documentos complejos con ecuaciones matemáticas, tablas multinivel, y referencias cruzadas automáticas sin pérdida de formato. Además, su código abierto bajo LPPL (LaTeX Project Public License) fomenta contribuciones comunitarias, con más de 5,000 paquetes en CTAN \cite{ieeetran_doc}.

\subsection{Extensión Masiva}

El repositorio CTAN alberga 6,907 paquetes (Abril 2026) que extienden LaTeX \cite{ctan_stats}:

\begin{itemize}
    \item \textbf{graphicx, tikz}: más de 2,500 paquetes gráficos
    \item \textbf{amsmath, mathtools}: más de 800 paquetes matemáticos
    \item \textbf{biblatex, natbib}: más de 400 gestores bibliográficos
    \item \textbf{beamer, tcolorbox}: más de 1,200 presentaciones/interfaces de usuario
\end{itemize}

\subsection{Importancia Estratégica}

\LaTeX presenta las siguientes ventajas estratégicas para investigadores y académicos:

\begin{enumerate}
    \item \textbf{Compilación reproducible}: se obtiene siempre el mismo PDF a partir del mismo código fuente, crucial para la integridad académica.
    \item \textbf{Ecuaciones semánticas}: e.g. $\nabla\cdot\mathbf{E}=\rho/\epsilon_0$ vs imágenes estáticas, mejor para accesibilidad y búsqueda.
    \item \textbf{Referencias cruzadas}: referencias automáticas a figuras, tablas, ecuaciones y secciones sin preocuparse por la numeración.
    \item \textbf{Bibliografía dinámica}: referencias a citas bibliográficas sin arduo trabajo manual.
    \item \textbf{Estilos journal oficiales}: e.g. IEEEtran.cls, Nature.cls, Springer.cls garantizan cumplimiento de formatos editoriales.
\end{enumerate}